\section{In Pestilence and Bereavement: Cyprian, Ambrose, and Augustine}

Doctrine about death is tested not in a quiet study, but when pestilence empties a city or a beloved hand grows cold. Three enduring monuments of Western patristic literature show Christian pastors speaking directly from the midst of epidemic and grief: St.~Cyprian of Carthage preaching through a devastating plague, St.~Ambrose of Milan standing beside his brother's bier, and St.~Augustine of Hippo remembering his mother Monica at Ostia.

\subsection{St.~Cyprian of Carthage in Time of Plague}

Around the year 252, a virulent pestilence swept the Roman Empire, striking Carthage with bowel flux, fever burning in the marrow, inflamed eyes, and rotting limbs. Watching his congregation shaken because baptized believers died side by side with pagans---and some grieved that disease might rob them of martyrdom---Bishop Cyprian wrote \emph{De mortalitate} (\emph{On the Mortality}).\endnote{St.~Cyprian of Carthage, \emph{On the Mortality} (\emph{De mortalitate}), \S\S~2--3, 8--18 (Ernest Wallis trans., ANF, vol.~5).}

He speaks with bracing candor. ``For it is for him to fear death who is not willing to go to Christ,'' Cyprian writes; ``it is for him to be unwilling to go to Christ who does not believe that he is about to reign with Christ'' (\S~2). Setting aged Simeon's \emph{Nunc dimittis} (Luke 2:29--30) before them, he shows that God's servants enter true peace only when withdrawn from this world's storms into ``the harbour of our home and eternal security'' (\S~3). Sharing common mortal clay with pagans does not make the plague the same event for both: ``The righteous are called to their place of refreshing, the unrighteous are snatched away to punishment'' (\S~15). Like Job and Tobias, Christians are proved in the tempest---``For the helmsman is recognised in the tempest; in the warfare the soldier is proved'' (\S~12)---while the epidemic tests whether the healthy tend the sick and the rich open their hands to the poor (\S~16). To those grieving that fever denies them red martyrdom, Cyprian answers tenderly: ``It is one thing for the spirit to be wanting for martyrdom, and another for martyrdom to have been wanting for the spirit\ldots{} For God does not ask for our blood, but for our faith'' (\S~17).

Exposing the contradiction of praying daily ``Thy kingdom come'' while recoiling when God calls us home (\S~18), Cyprian reminds his flock that the departed ``are not lost, but sent before'' (\emph{non amitti sed praemitti}), and that bodily death is ``not an ending, but a transit, and, this journey of time being traversed, a passage to eternity.''\endnote{St.~Cyprian, \emph{On the Mortality}, \S\S~20, 22, 26 (ANF, vol.~5).} Then, in his closing chapter, he lifts their eyes above the plague-pits to their true fatherland:

\begin{quote}
We regard paradise as our country---we already begin to consider the patriarchs as our parents: why do we not hasten and run, that we may behold our country, that we may greet our parents? There a great number of our dear ones is awaiting us, and a dense crowd of parents, brothers, children, is longing for us, already assured of their own safety, and still solicitous for our salvation. To attain to their presence and their embrace, what a gladness both for them and for us in common! [\ldots] There the glorious company of the apostles---there the host of the rejoicing prophets---there the innumerable multitude of martyrs, crowned for the victory of their struggle and passion. (\S~26)
\end{quote}

\subsection{St.~Ambrose Beside the Bier of Satyrus}

A century and a quarter later, St.~Ambrose of Milan buried his elder brother Satyrus, who had surrendered public life to manage Ambrose's household after his election as bishop and died in his arms after surviving shipwreck on a voyage to Africa.\endnote{St.~Ambrose of Milan, \emph{On the Decease of His Brother Satyrus} (\emph{De excessu fratris Satyri}), Books I--II (H.~de Romestin trans., NPNF, ser.~2, vol.~10, 1896). Book I was preached at the funeral; Book II seven days later at the grave.} Standing beside his brother's body, Ambrose never pretends that faith numbs pain. He speaks of trying to catch Satyrus's last breath in a parting kiss, comparing himself to an ox lowing after its missing yoke-fellow (I.7--8, 19). Yet he gathers that anguish immediately into sacrifice and gratitude: ``We have brought hither, dearest brethren, my sacrifice, a sacrifice undefiled, a sacrifice well pleasing to God, my lord and brother Satyrus\ldots{} I must rather rejoice that I had such a brother than grieve that I had lost a brother, for the former is a gift, the latter a debt to be paid'' (I.1, 3--4).

Above all, Ambrose vindicates the dignity of Christian tears, distinguishing natural sorrow from faithless despair:

\begin{quote}
But we have not incurred any grievous sin by our tears. Not all weeping proceeds from unbelief or weakness. Natural grief is one thing, distrustful sadness is another, and there is a very great difference between longing for what you have lost and lamenting that you have lost it. [\ldots] I confess, then, that I too wept, but the Lord also wept. He wept for one not related to Him, I for my brother. He wept for all in weeping for one, I will weep for thee in all, my brother. (I.10; cf.\ John 11:35)
\end{quote}

``I have not lost but changed my intercourse with thee,'' Ambrose tells Satyrus; ``thou art become surety for that other country, for I begin to be no stranger there where the better portion of myself already is'' (I.6). Returning to the tomb on the seventh day for Book II, Ambrose distinguishes spiritual death to sin, the dreadful death of guilt, and bodily departure (II.36--37), teaching that in a fallen world God's mercy made bodily death a \emph{remedy} setting a limit to earthly misery:

\begin{quote}
And, indeed, death was no part of man's nature, but became natural; for God did not institute death at first, but gave it as a remedy. [\ldots] Because of transgression the life of man, condemned to lengthened labour, began to be wretched with intolerable groaning; so that it was fitting that an end should be set to the evils, and that death should restore what life had lost. For immortality, unless grace breathed upon it, would be rather a burden than an advantage. And if one consider accurately, it is not the death of our being, but of evil, for being continues, it is evil that perishes. (II.47--48; cf.\ II.38)\endnote{St.~Ambrose, \emph{On the Decease of His Brother Satyrus}, II.46--48 (NPNF2 vol.~10).}
\end{quote}

In Christ, that boundary becomes victory: ``when we offer the Sacrifice we declare His death, for His death is victory, His death is our mystery, His death is the yearly recurring solemnity of the world'' (II.46).

\subsection{St.~Augustine and St.~Monica at Ostia}

Ten years later, Ambrose's spiritual son Augustine stood with his mother Monica at a garden window in Ostia, waiting for passage back to Africa after his Easter Baptism in 387. Speaking together of the eternal life of the saints, their minds rose above sun, stars, and changing speech until for a single instant they touched Eternal Wisdom itself, leaving there ``the first-fruits of the Spirit.''\endnote{St.~Augustine, \emph{Confessions}, IX.10.23--26 (J.~G.~Pilkington trans., NPNF, ser.~1, vol.~1).} In the stillness afterward, Monica spoke her own \emph{Nunc dimittis}:

\begin{quote}
Son, for myself, I have no longer any pleasure in aught in this life. What I want here further, and why I am here, I know not, now that my hopes in this world are satisfied. There was indeed one thing for which I wished to tarry a little in this life, and that was that I might see thee a Catholic Christian before I died. My God has exceeded this abundantly, so that I see thee despising all earthly felicity, made His servant,---what do I here? (IX.10.26)
\end{quote}

Five days later, fever struck. Waking from a swoon, she told her sons calmly, ``Here shall you bury your mother.'' Once she had carefully prepared a grave beside her husband Patricius in Africa; now grace had detached her even from that care: ``Nothing is far to God; nor need I fear lest He should be ignorant at the end of the world of the place whence He is to raise me up.''\endnote{St.~Augustine, \emph{Confessions}, IX.11.27--28 (NPNF1 vol.~1).} To her sons she left one enduring testament:

\begin{quote}
Lay this body anywhere, let not the care for it trouble you at all. This only I ask, that you will remember me at the Lord's altar, wherever you be. (IX.11.27)
\end{quote}

On the ninth day of her illness, aged fifty-six, she died. Augustine fought back his tears at the funeral while Evodius sang Psalm 100 (``Mercy and judgment I will sing to thee, O Lord'') and ``the sacrifice of our redemption'' was offered beside the grave. Inside, however, his heart was torn by ``that newly-made wound, from having that most sweet and dear habit of living together suddenly broken off,'' compounded by shame over his own sorrow. Only after waking with Ambrose's evening hymn \emph{Deus creator omnium} upon his lips did the hard crust yield, and alone before God he let the tears flow freely, ``spreading them beneath my heart; and it rested in them, for Thy ears were nigh me.''\endnote{St.~Augustine, \emph{Confessions}, IX.12.29--13.37 (NPNF1 vol.~1).} Even then, he poured out ``tears of a far different sort,'' begging pardon for whatever frailties Monica had contracted since Baptism, appealing to ``that Medicine of our wounds who hung upon the tree,'' and asking every reader to remember Monica and Patricius at the altar of God.

In Cyprian's longing for the heavenly fatherland, Ambrose's defense of tears hallowed by Christ, and Augustine's altar-bound vigil beside Monica, the early Church displays a wholeness that neither denies human love nor bows to despair.
